\documentclass[11pt]{article}

\usepackage[margin=1in]{geometry}
\usepackage{amsmath,amssymb,amsthm}
\usepackage{mathrsfs}
\usepackage{graphicx}
\usepackage[hidelinks]{hyperref}

\newtheorem{theorem}{Theorem}
\newtheorem{lemma}{Lemma}

\newcommand{\C}{\mathbb C}
\newcommand{\T}{\mathbb T}
\newcommand{\FG}{\mathscr F_G}
\newcommand{\ellone}{\ell_1^2}
\newcommand{\de}{\delta_e}
\newcommand{\ds}{\delta_s}
\newcommand{\ze}{\zeta_e}
\newcommand{\zs}{\zeta_s}

\title{A Two-Point Group Frame Whose Generator Set Is Disconnected}
\author{arXiv automatic math run \texttt{fa\_banach\_001}}
\date{June 28, 2026}

\begin{document}
\maketitle

\section*{Status}

This packet gives a candidate counterexample to Problem 2.19 of Krishna
\cite{Krishna2023}. The problem asks whether the set
\(\FG(\pi,\mathcal X)\subseteq \mathcal X^*\times\mathcal X\) of group-p-USF
functional-vectors is path connected in the norm topology. The answer is no
already for \(p=1\), the two-element group, and the left regular representation
on \(\ell_1^2(\C)\).

\section*{Source Question}

The source asks the following at the end of Section 2.

\begin{center}
\includegraphics[width=\textwidth]{figures/open_problem_crop.png}
\end{center}

We address Problem 2.19. The same calculation also gives a complete
description of \(\FG(\pi,\mathcal X)\) for this particular representation,
which is a small instance of Problem 2.18.

\section*{Definitions Used}

Let \(G\) be a discrete group and let \(\pi:G\to \mathbb I^*(\mathcal X)\) be
an invertible isometric representation. Given \(f\in\mathcal X^*\) and
\(\tau\in\mathcal X\), the associated group-generated functionals and vectors
are
\[
  f_g=f\pi_{g^{-1}},\qquad \tau_g=\pi_g\tau.
\]
The pair \((f,\tau)\) belongs to \(\FG(\pi,\mathcal X)\) when
\((\{f_g\}_{g\in G},\{\tau_g\}_{g\in G})\) is a group-p-USF. For \(p=1\) and
finite \(G\), this means that the analysis map
\[
  \theta_f x=(f_g(x))_{g\in G}:\mathcal X\to \ell_1(G)
\]
is an isometry and the synthesis map
\[
  \theta_\tau(a_g)_{g\in G}=\sum_{g\in G}a_g\tau_g
\]
is bounded and satisfies the unconditional Schauder reconstruction
\(\theta_\tau\theta_f=I_{\mathcal X}\).

\section*{The Counterexample}

Let \(G=C_2=\{e,s\}\), with \(s^2=e\). Put
\(\mathcal X=\ell_1(G)=\ell_1^2(\C)\), with standard basis
\(\de,\ds\) and coordinate functionals \(\ze,\zs\). Let \(\pi=\lambda\) be the
left regular representation:
\[
  \lambda_s\de=\ds,\qquad \lambda_s\ds=\de.
\]
Then \(\lambda\) is an invertible isometric representation on \(\ell_1^2(\C)\).
The pair \((\ze,\de)\) is a group-1-USF functional-vector because its analysis
and synthesis maps are both the identity on \(\ell_1^2\). Hence
\(\FG(\lambda,\ell_1^2)\neq\varnothing\).

\begin{lemma}
With the preceding notation,
\[
\FG(\lambda,\ell_1^2)
=
\{(\alpha\ze,\alpha^{-1}\de):\alpha\in\T\}
\cup
\{(\alpha\zs,\alpha^{-1}\ds):\alpha\in\T\}.
\]
\end{lemma}

\begin{proof}
Take any \((f,\tau)\in\FG(\lambda,\ell_1^2)\). Write
\[
  f=a\ze+b\zs,\qquad \tau=u\de+v\ds.
\]
Relative to the ordered basis \((\de,\ds)\) of both the domain and
\(\ell_1(G)\), the analysis operator is
\[
  A=\theta_f=
  \begin{pmatrix}
    a & b\\
    b & a
  \end{pmatrix},
\]
because the second coordinate functional is
\(f_s=f\lambda_s=b\ze+a\zs\). The group-1-USF condition says that
\(A:\ell_1^2\to\ell_1^2\) is an isometry. Since the spaces have the same
finite dimension, \(A\) is a surjective linear isometry.

Surjective linear isometries of complex \(\ell_1^2\) are monomial. Indeed, an
isometry maps the extreme points of the unit ball onto extreme points, and the
extreme points of the complex \(\ell_1^2\)-unit ball are precisely
\(\gamma\de\) and \(\gamma\ds\), \(|\gamma|=1\). Therefore each column of
\(A\) has exactly one nonzero coordinate of modulus one, and the two columns
have different supports.

The symmetry of \(A\) leaves only two possibilities:
\[
  A=\alpha I\quad(|\alpha|=1),
  \qquad\text{or}\qquad
  A=\alpha
  \begin{pmatrix}
  0&1\\
  1&0
  \end{pmatrix}
  \quad(|\alpha|=1).
\]

The synthesis operator determined by \(\tau=u\de+v\ds\) has matrix
\[
  B=\theta_\tau=
  \begin{pmatrix}
    u & v\\
    v & u
  \end{pmatrix},
\]
because its columns are \(\tau_e=\tau\) and \(\tau_s=\lambda_s\tau\). The USF
reconstruction gives \(BA=I\). Hence \(B=A^{-1}\).

If \(A=\alpha I\), then \(B=\alpha^{-1}I\), so
\((f,\tau)=(\alpha\ze,\alpha^{-1}\de)\). If \(A=\alpha S\), where \(S\) is the
swap matrix, then \(B=\alpha^{-1}S\), so
\((f,\tau)=(\alpha\zs,\alpha^{-1}\ds)\). Conversely, each pair in these two
families plainly has \(A=\alpha I\) or \(A=\alpha S\) and
\(B=A^{-1}\), so it is a group-1-USF functional-vector.
\end{proof}

\begin{theorem}
\(\FG(\lambda,\ell_1^2(\C))\subseteq(\ell_1^2)^*\times\ell_1^2\) is not path
connected in the norm topology.
\end{theorem}

\begin{proof}
By the lemma, \(\FG(\lambda,\ell_1^2)\) is the union of the two nonempty sets
\[
  E_e=\{(\alpha\ze,\alpha^{-1}\de):\alpha\in\T\},
  \qquad
  E_s=\{(\alpha\zs,\alpha^{-1}\ds):\alpha\in\T\}.
\]
These sets are separated. For any \(\alpha,\beta\in\T\),
\[
  \|\alpha\ze-\beta\zs\|_{(\ell_1^2)^*}
  =
  \|(\alpha,-\beta)\|_{\ell_\infty^2}
  =1.
\]
Thus, in any of the standard equivalent product norms on
\((\ell_1^2)^*\times\ell_1^2\), the distance between \(E_e\) and \(E_s\) is
positive. Consequently \(E_e\) and \(E_s\) are disjoint relatively open and
closed subsets of \(\FG(\lambda,\ell_1^2)\). Since both are nonempty, the set
\(\FG(\lambda,\ell_1^2)\) is disconnected, and hence cannot be path connected.
\end{proof}

\section*{Verification Notes}

No infinite-dimensional or convergence issue is hidden in the construction:
the group is finite and all sums have two terms. The only structural input is
the elementary finite-dimensional classification of the complex
\(\ell_1^2\)-isometries. The proof also shows that, for this representation,
the inclusions from Proposition 2.14 and Theorem 2.17 are equalities:
the invertible isometries commuting with the swap are exactly the scalar
identity and scalar swap operators.

\section*{Novelty Check}

The local run indexes were searched for arXiv:2305.01499, the title
\emph{Group-Frames for Banach Spaces}, group-p-USF, Banach frames, frame,
group representation, orbit, Schauder frame, path connected, and
\(\mathscr F_G\) keywords. No exact prior packet for this question was present.

External searches on June 28, 2026 for combinations of arXiv:2305.01499,
the source title, path connected, group-p-USF, and \(\mathscr F_G\) did not
locate a published answer. This should be treated as a bounded novelty check.

\section*{Human Review Recommendation}

Confirm that Problem 2.19 is meant for the fixed \(p\) in the surrounding
group-p-USF context. Under that reading, choosing \(p=1\) gives a complete
negative answer to the general path-connectedness question.

\begin{thebibliography}{9}

\bibitem{Krishna2023}
K. Mahesh Krishna,
\emph{Group-Frames for Banach Spaces},
arXiv:2305.01499.

\end{thebibliography}

\end{document}
